% Options for packages loaded elsewhere
\PassOptionsToPackage{unicode}{hyperref}
\PassOptionsToPackage{hyphens}{url}
\documentclass[
]{article}
\usepackage{xcolor}
\usepackage[margin=1in]{geometry}
\usepackage{amsmath,amssymb}
\setcounter{secnumdepth}{-\maxdimen} % remove section numbering
\usepackage{iftex}
\ifPDFTeX
  \usepackage[T1]{fontenc}
  \usepackage[utf8]{inputenc}
  \usepackage{textcomp} % provide euro and other symbols
\else % if luatex or xetex
  \usepackage{unicode-math} % this also loads fontspec
  \defaultfontfeatures{Scale=MatchLowercase}
  \defaultfontfeatures[\rmfamily]{Ligatures=TeX,Scale=1}
\fi
\usepackage{lmodern}
\ifPDFTeX\else
  % xetex/luatex font selection
\fi
% Use upquote if available, for straight quotes in verbatim environments
\IfFileExists{upquote.sty}{\usepackage{upquote}}{}
\IfFileExists{microtype.sty}{% use microtype if available
  \usepackage[]{microtype}
  \UseMicrotypeSet[protrusion]{basicmath} % disable protrusion for tt fonts
}{}
\makeatletter
\@ifundefined{KOMAClassName}{% if non-KOMA class
  \IfFileExists{parskip.sty}{%
    \usepackage{parskip}
  }{% else
    \setlength{\parindent}{0pt}
    \setlength{\parskip}{6pt plus 2pt minus 1pt}}
}{% if KOMA class
  \KOMAoptions{parskip=half}}
\makeatother
\usepackage{color}
\usepackage{fancyvrb}
\newcommand{\VerbBar}{|}
\newcommand{\VERB}{\Verb[commandchars=\\\{\}]}
\DefineVerbatimEnvironment{Highlighting}{Verbatim}{commandchars=\\\{\}}
% Add ',fontsize=\small' for more characters per line
\usepackage{framed}
\definecolor{shadecolor}{RGB}{248,248,248}
\newenvironment{Shaded}{\begin{snugshade}}{\end{snugshade}}
\newcommand{\AlertTok}[1]{\textcolor[rgb]{0.94,0.16,0.16}{#1}}
\newcommand{\AnnotationTok}[1]{\textcolor[rgb]{0.56,0.35,0.01}{\textbf{\textit{#1}}}}
\newcommand{\AttributeTok}[1]{\textcolor[rgb]{0.13,0.29,0.53}{#1}}
\newcommand{\BaseNTok}[1]{\textcolor[rgb]{0.00,0.00,0.81}{#1}}
\newcommand{\BuiltInTok}[1]{#1}
\newcommand{\CharTok}[1]{\textcolor[rgb]{0.31,0.60,0.02}{#1}}
\newcommand{\CommentTok}[1]{\textcolor[rgb]{0.56,0.35,0.01}{\textit{#1}}}
\newcommand{\CommentVarTok}[1]{\textcolor[rgb]{0.56,0.35,0.01}{\textbf{\textit{#1}}}}
\newcommand{\ConstantTok}[1]{\textcolor[rgb]{0.56,0.35,0.01}{#1}}
\newcommand{\ControlFlowTok}[1]{\textcolor[rgb]{0.13,0.29,0.53}{\textbf{#1}}}
\newcommand{\DataTypeTok}[1]{\textcolor[rgb]{0.13,0.29,0.53}{#1}}
\newcommand{\DecValTok}[1]{\textcolor[rgb]{0.00,0.00,0.81}{#1}}
\newcommand{\DocumentationTok}[1]{\textcolor[rgb]{0.56,0.35,0.01}{\textbf{\textit{#1}}}}
\newcommand{\ErrorTok}[1]{\textcolor[rgb]{0.64,0.00,0.00}{\textbf{#1}}}
\newcommand{\ExtensionTok}[1]{#1}
\newcommand{\FloatTok}[1]{\textcolor[rgb]{0.00,0.00,0.81}{#1}}
\newcommand{\FunctionTok}[1]{\textcolor[rgb]{0.13,0.29,0.53}{\textbf{#1}}}
\newcommand{\ImportTok}[1]{#1}
\newcommand{\InformationTok}[1]{\textcolor[rgb]{0.56,0.35,0.01}{\textbf{\textit{#1}}}}
\newcommand{\KeywordTok}[1]{\textcolor[rgb]{0.13,0.29,0.53}{\textbf{#1}}}
\newcommand{\NormalTok}[1]{#1}
\newcommand{\OperatorTok}[1]{\textcolor[rgb]{0.81,0.36,0.00}{\textbf{#1}}}
\newcommand{\OtherTok}[1]{\textcolor[rgb]{0.56,0.35,0.01}{#1}}
\newcommand{\PreprocessorTok}[1]{\textcolor[rgb]{0.56,0.35,0.01}{\textit{#1}}}
\newcommand{\RegionMarkerTok}[1]{#1}
\newcommand{\SpecialCharTok}[1]{\textcolor[rgb]{0.81,0.36,0.00}{\textbf{#1}}}
\newcommand{\SpecialStringTok}[1]{\textcolor[rgb]{0.31,0.60,0.02}{#1}}
\newcommand{\StringTok}[1]{\textcolor[rgb]{0.31,0.60,0.02}{#1}}
\newcommand{\VariableTok}[1]{\textcolor[rgb]{0.00,0.00,0.00}{#1}}
\newcommand{\VerbatimStringTok}[1]{\textcolor[rgb]{0.31,0.60,0.02}{#1}}
\newcommand{\WarningTok}[1]{\textcolor[rgb]{0.56,0.35,0.01}{\textbf{\textit{#1}}}}
\usepackage{longtable,booktabs,array}
\usepackage{calc} % for calculating minipage widths
% Correct order of tables after \paragraph or \subparagraph
\usepackage{etoolbox}
\makeatletter
\patchcmd\longtable{\par}{\if@noskipsec\mbox{}\fi\par}{}{}
\makeatother
% Allow footnotes in longtable head/foot
\IfFileExists{footnotehyper.sty}{\usepackage{footnotehyper}}{\usepackage{footnote}}
\makesavenoteenv{longtable}
\usepackage{graphicx}
\makeatletter
\newsavebox\pandoc@box
\newcommand*\pandocbounded[1]{% scales image to fit in text height/width
  \sbox\pandoc@box{#1}%
  \Gscale@div\@tempa{\textheight}{\dimexpr\ht\pandoc@box+\dp\pandoc@box\relax}%
  \Gscale@div\@tempb{\linewidth}{\wd\pandoc@box}%
  \ifdim\@tempb\p@<\@tempa\p@\let\@tempa\@tempb\fi% select the smaller of both
  \ifdim\@tempa\p@<\p@\scalebox{\@tempa}{\usebox\pandoc@box}%
  \else\usebox{\pandoc@box}%
  \fi%
}
% Set default figure placement to htbp
\def\fps@figure{htbp}
\makeatother
\setlength{\emergencystretch}{3em} % prevent overfull lines
\providecommand{\tightlist}{%
  \setlength{\itemsep}{0pt}\setlength{\parskip}{0pt}}
\usepackage{bookmark}
\IfFileExists{xurl.sty}{\usepackage{xurl}}{} % add URL line breaks if available
\urlstyle{same}
\hypersetup{
  hidelinks,
  pdfcreator={LaTeX via pandoc}}

\author{}
\date{\vspace{-2.5em}}

\begin{document}

\section{Réseaux mobiles}\label{ruxe9seaux-mobiles}

\subsection{Objectif de cet ouvrage}\label{objectif-de-cet-ouvrage}

L'objectif de ce livre est d'expliquer le fonctionnement interne des
réseaux mobiles en partant des bases radio jusqu'aux mécanismes de
signalisation qui orchestrent leur fonctionnement.

À travers l'étude d'un laboratoire reproductible, nous analyserons :

\begin{itemize}
\tightlist
\item
  l'architecture des réseaux GSM
\item
  les protocoles de signalisation utilisés par les opérateurs
\item
  les interactions entre réseaux mobiles
\item
  les mécanismes historiques qui ont façonné l'infrastructure actuelle
\end{itemize}

Cette approche permet de comprendre non seulement comment ces réseaux
fonctionnent, mais aussi pourquoi ils ont été conçus de cette manière.

Les chapitres suivants exploreront progressivement les différentes
couches du système, en commençant par les bases du GSM et de la
signalisation SS7, avant d'aborder les architectures plus modernes et
les implications en matière de sécurité.

\subsubsection{Architecture GSM/EGPRS}\label{architecture-gsmegprs}

\begin{center}\includegraphics[width=4.35in]{../../../../tmp/RtmpVUNkQ4/file28a5522632abc} \end{center}

\subsection{Présentation}\label{pruxe9sentation}

Ce guide décrit l'installation modulaire d'une infrastructure
\textbf{EGPRS} (GPRS/EDGE) basée sur Osmocom et SDR (USRP, BladeRF).\\
Chaque chapitre détaille un module, son rôle, ses dépendances et
options.\\
Tout peut être scripté pour une intégration Dockerfile ou un déploiement
manuel sur Debian/Ubuntu 22.04+.

\begin{center}\rule{0.5\linewidth}{0.5pt}\end{center}

\subsection{Dépendances système}\label{duxe9pendances-systuxe8me}

Commencez par installer tous les outils de build essentiels,
bibliothèques, Python et utilitaires de debug : - Compilation :
\texttt{gcc}, \texttt{g++}, \texttt{make}, \texttt{cmake},
\texttt{autoconf}, \texttt{automake}, \texttt{libtool}, etc. - Libs
système : \texttt{libusb-1.0-0-dev}, \texttt{libedit-dev},
\texttt{libfftw3-dev}, \texttt{liburing-dev}, \ldots{} - Dev Python :
\texttt{python3-dev}, \texttt{python3-pip}, \texttt{python3-numpy},
\texttt{python3-mako}, \texttt{python3-lxml}, \texttt{python3-six},
\ldots{} - Outils additionnels : \texttt{nano}, \texttt{git},
\texttt{wget}, \texttt{curl}, \texttt{tmux}, \texttt{doxygen},
\texttt{ethtool}, etc.

\textbf{Commande d'installation :}

\begin{Shaded}
\begin{Highlighting}[]
\ExtensionTok{apt}\NormalTok{ update}
\ExtensionTok{apt}\NormalTok{ install }\AttributeTok{{-}y} \OperatorTok{\textless{}}\NormalTok{tous{-}les{-}paquets{-}listés}\OperatorTok{\textgreater{}}
\end{Highlighting}
\end{Shaded}

\begin{center}\rule{0.5\linewidth}{0.5pt}\end{center}

\subsection{OsmocomBB}\label{osmocombb}

\textbf{OsmocomBB} (Osmocom Baseband) est un projet open source qui
implémente une pile GSM ``baseband'' pour certains téléphones Motorola
et des configurations SDR.

\subsubsection{Qu'est-ce qu'OsmocomBB\,?}\label{quest-ce-quosmocombb}

\begin{itemize}
\tightlist
\item
  \textbf{Pile protocolaire GSM côté mobile (MS, mobile station)}
\item
  Remplace le firmware d'origine pour donner accès aux couches GSM
  1/2/3, journalisation des protocoles, sniffing et tests
\item
  Utilisations\,: analyse, tests de sécurité, recherche académique, IMSI
  catchers, fuzzing radio
\end{itemize}

\subsubsection{Architecture}\label{architecture}

\begin{itemize}
\tightlist
\item
  \textbf{host/osmocon}\,: Outil de chargement du firmware sur les
  appareils compatibles
\item
  \textbf{src/host/layer23}\,: Pile GSM s'exécutant sur le PC (appels,
  SMS, sniffing, capture de bursts, etc.)
\item
  \textbf{firmware/}\,: Code baseband à flasher sur les téléphones
  compatibles
\end{itemize}

\begin{center}\includegraphics[width=1.16in]{../../../../tmp/RtmpVUNkQ4/file28a5562e94f83} \end{center}

\subsubsection{Matériel supporté}\label{matuxe9riel-supportuxe9}

\begin{itemize}
\tightlist
\item
  \textbf{Série Motorola C1xx (C123, C118, C155, etc.)}
\item
  Certains autres modèles à chipset ``Calypso''
\end{itemize}

\subsubsection{Cas d'usage \& limitations}\label{cas-dusage-limitations}

\begin{itemize}
\tightlist
\item
  Peut se connecter à une vraie cellule 2G ou à une BTS Osmocom
  (osmo-bts-trx)
\item
  Limité à la 2G pure (pas de support 3G/4G)
\item
  \textbf{Attention\,:} Émettre sans licence est illégal\,!\\
  N'utilisez que pour la recherche ou en environnement contrôlé.
\end{itemize}

\subsubsection{Liens utiles}\label{liens-utiles}

\begin{itemize}
\tightlist
\item
  \href{https://github.com/osmocom/osmocom-bb}{Dépôt GitHub OsmocomBB}
\item
  \href{https://osmocom.org/projects/baseband/wiki/OsmocomBB}{Wiki
  officiel}
\item
  \href{https://www.lafabriquedunet.fr/blog/osmocombb-guide/}{Guide
  d'installation (FR)}
\end{itemize}

\begin{center}\rule{0.5\linewidth}{0.5pt}\end{center}

\textbf{Pour SDR\,:} Un projet expérimental existe pour connecter
OsmocomBB à un SDR (BladeRF, HackRF, etc.) via l'interface ``virtual
layer1''.

\subsection{OsmocomBB Mobile (MS) -- Installation rapide Ubuntu
24.04}\label{osmocombb-mobile-ms-installation-rapide-ubuntu-24.04}

\subsubsection{Prérequis}\label{pruxe9requis}

\begin{itemize}
\item
  \textbf{Matériel\,:}

  \begin{itemize}
  \tightlist
  \item
    Motorola C1xx (C115/C118/C123/C155, etc., baseband Calypso)
  \item
    Câble USB-série 2.5\,mm (PL2303 ou FTDI)
  \end{itemize}
\item
  \textbf{Système\,:}

  \begin{itemize}
  \tightlist
  \item
    Ubuntu 24.04 x86\_64 (natif ou VM)
  \item
    Accès internet + droits sudo/root
  \end{itemize}
\end{itemize}

\begin{center}\rule{0.5\linewidth}{0.5pt}\end{center}

\subsubsection{Installer les
dépendances}\label{installer-les-duxe9pendances}

\begin{Shaded}
\begin{Highlighting}[]
\FunctionTok{sudo}\NormalTok{ apt update}
\FunctionTok{sudo}\NormalTok{ apt install }\AttributeTok{{-}y} \DataTypeTok{\textbackslash{}}
\NormalTok{  build{-}essential git cmake pkg{-}config autoconf automake libtool shtool }\DataTypeTok{\textbackslash{}}
\NormalTok{  bison flex libtalloc{-}dev libpcsclite{-}dev zlib1g{-}dev libmpfr{-}dev libmpc{-}dev }\DataTypeTok{\textbackslash{}}
\NormalTok{  libgmp{-}dev libncurses5{-}dev libncursesw5{-}dev libsofia{-}sip{-}ua{-}dev libxml2{-}dev }\DataTypeTok{\textbackslash{}}
\NormalTok{  libfftw3{-}dev libgnutls28{-}dev libssl{-}dev python3 python3{-}pip }\DataTypeTok{\textbackslash{}}
\NormalTok{  alsa{-}oss lemon libtinfo{-}dev}
\end{Highlighting}
\end{Shaded}

\begin{center}\rule{0.5\linewidth}{0.5pt}\end{center}

\subsubsection{Installer la toolchain ARM (cross-compile
Calypso)}\label{installer-la-toolchain-arm-cross-compile-calypso}

\begin{Shaded}
\begin{Highlighting}[]
\BuiltInTok{cd}\NormalTok{ /opt}
\FunctionTok{sudo}\NormalTok{ wget https://downloads.osmocom.org/binaries/toolchain/arm{-}2021.03.tar.bz2}
\FunctionTok{sudo}\NormalTok{ tar }\AttributeTok{{-}xjf}\NormalTok{ arm{-}2021.03.tar.bz2}
\BuiltInTok{export} \VariableTok{PATH}\OperatorTok{=}\VariableTok{$PATH}\NormalTok{:/opt/arm{-}2021.03/bin}
\end{Highlighting}
\end{Shaded}

\begin{quote}
\emph{(Ajoute ce dernier export à ton \texttt{\textasciitilde{}/.bashrc}
pour qu'il soit permanent)}
\end{quote}

\begin{center}\rule{0.5\linewidth}{0.5pt}\end{center}

\subsubsection{Cloner et compiler
OsmocomBB}\label{cloner-et-compiler-osmocombb}

\begin{Shaded}
\begin{Highlighting}[]
\FunctionTok{git}\NormalTok{ clone https://github.com/osmocom/osmocom{-}bb.git}
\BuiltInTok{cd}\NormalTok{ osmocom{-}bb/src}
\ExtensionTok{autoreconf} \AttributeTok{{-}i}
\BuiltInTok{export} \VariableTok{PATH}\OperatorTok{=}\VariableTok{$PATH}\NormalTok{:/opt/arm{-}2021.03/bin}
\ExtensionTok{./configure}
\FunctionTok{make} \AttributeTok{{-}j}\VariableTok{$(}\FunctionTok{nproc}\VariableTok{)}
\end{Highlighting}
\end{Shaded}

\begin{center}\rule{0.5\linewidth}{0.5pt}\end{center}

\subsubsection{Flasher le firmware Mobile (Compal/C123
etc.)}\label{flasher-le-firmware-mobile-compalc123-etc.}

\begin{enumerate}
\def\labelenumi{\arabic{enumi}.}
\item
  \textbf{Branche le téléphone via le câble USB-série (sans la
  batterie).}
\item
  \textbf{Dans un terminal\,:}
\end{enumerate}

\begin{Shaded}
\begin{Highlighting}[]
\BuiltInTok{cd}\NormalTok{ host/osmocon}
\FunctionTok{sudo}\NormalTok{ ./osmocon }\AttributeTok{{-}m}\NormalTok{ c123xor }\AttributeTok{{-}p}\NormalTok{ /dev/ttyUSB0 }\AttributeTok{{-}c}\NormalTok{ ../../target/firmware/board/compal\_e88/hello\_world.highram.bin}
\CommentTok{\# (optionnel : tester hello\_world)}
\end{Highlighting}
\end{Shaded}

\begin{enumerate}
\def\labelenumi{\arabic{enumi}.}
\setcounter{enumi}{2}
\tightlist
\item
  \textbf{Pour le firmware mobile complet\,:}
\end{enumerate}

\begin{Shaded}
\begin{Highlighting}[]
\FunctionTok{sudo}\NormalTok{ ./osmocon }\AttributeTok{{-}m}\NormalTok{ c123xor }\AttributeTok{{-}p}\NormalTok{ /dev/ttyUSB0 }\AttributeTok{{-}c}\NormalTok{ ../../target/firmware/board/compal\_e88/layer1.highram.bin}
\end{Highlighting}
\end{Shaded}

\begin{enumerate}
\def\labelenumi{\arabic{enumi}.}
\setcounter{enumi}{3}
\tightlist
\item
  \textbf{Dans un autre terminal\,:}
\end{enumerate}

\begin{Shaded}
\begin{Highlighting}[]
\BuiltInTok{cd}\NormalTok{ host/layer23/src/mobile}
\ExtensionTok{./mobile}
\end{Highlighting}
\end{Shaded}

\begin{center}\rule{0.5\linewidth}{0.5pt}\end{center}

\subsubsection{\texorpdfstring{Configurer \texttt{mobile.cfg} (SIM,
PLMN, log,
etc.)}{Configurer mobile.cfg (SIM, PLMN, log, etc.)}}\label{configurer-mobile.cfg-sim-plmn-log-etc.}

Crée/édite le fichier \texttt{mobile.cfg} (exemple dans
\texttt{osmocom-bb/src/host/layer23/src/mobile/})\,:

\begin{Shaded}
\begin{Highlighting}[]
\DataTypeTok{log gsmtap 127.0.0.1}
\DataTypeTok{log stderr}
\DataTypeTok{hlr}
\DataTypeTok{  subscriber{-}list subscriber\_list.csv}
\DataTypeTok{sim}
\DataTypeTok{  driver sysmo}
\end{Highlighting}
\end{Shaded}

\emph{(Adapte la section SIM si tu utilises une vraie SIM ou une ``test
SIM'')}

\begin{center}\rule{0.5\linewidth}{0.5pt}\end{center}

\subsubsection{Lancer OsmocomBB Mobile}\label{lancer-osmocombb-mobile}

\begin{Shaded}
\begin{Highlighting}[]
\CommentTok{\# Dans host/layer23/src/mobile/}
\ExtensionTok{./mobile} \AttributeTok{{-}c}\NormalTok{ mobile.cfg}
\end{Highlighting}
\end{Shaded}

\begin{itemize}
\tightlist
\item
  Le téléphone démarre, tente de s'enregistrer sur le réseau configuré.
\item
  Regarde les logs\,: tu dois voir ``PLMN search'', ``network found'',
  ``registration accepted'' (si SIM \& PLMN ok).
\end{itemize}

\begin{center}\rule{0.5\linewidth}{0.5pt}\end{center}

\subsubsection{(Optionnel) Suivi GSMTAP dans
Wireshark}\label{optionnel-suivi-gsmtap-dans-wireshark}

Wireshark peut décoder en temps réel les trames GSMTAP envoyées en
UDP\,:

\begin{Shaded}
\begin{Highlighting}[]
\ExtensionTok{wireshark} \AttributeTok{{-}k} \AttributeTok{{-}Y}\NormalTok{ gsmtap }\AttributeTok{{-}i}\NormalTok{ udp:4729}
\end{Highlighting}
\end{Shaded}

\emph{(Pense à bien avoir \texttt{log\ gsmtap\ 127.0.0.1} dans ta config
!)}

\begin{center}\rule{0.5\linewidth}{0.5pt}\end{center}

\subsubsection{Liens utiles}\label{liens-utiles-1}

\begin{itemize}
\tightlist
\item
  \href{https://osmocom.org/projects/baseband/wiki}{OsmocomBB Wiki}
\item
  \href{https://lists.osmocom.org/mailman/listinfo/baseband-devel}{Forum
  OsmocomBB}
\item
  \href{https://downloads.osmocom.org/binaries/toolchain/}{Firmware
  Toolchain}
\item
  \href{https://osmocom.org/projects/baseband/wiki/CalypsoBTS}{BTS
  HowTo}
\end{itemize}

\subsection{SDR \& Matériel (UHD,
BladeRF)}\label{sdr-matuxe9riel-uhd-bladerf}

\subsubsection{UHD (USRP/Ettus)}\label{uhd-usrpettus}

\begin{itemize}
\item
  \textbf{libuhd-dev, uhd-host} : Drivers et outils pour USRP.
\item
  \textbf{Images firmware} :
\end{itemize}

\begin{Shaded}
\begin{Highlighting}[]
\ExtensionTok{python3}\NormalTok{ /usr/lib/uhd/utils/uhd\_images\_downloader.py}
\end{Highlighting}
\end{Shaded}

\begin{itemize}
\tightlist
\item
  \textbf{Règle udev} pour détection auto :
\end{itemize}

\begin{Shaded}
\begin{Highlighting}[]
\FunctionTok{wget}\NormalTok{ https://raw.githubusercontent.com/EttusResearch/uhd/refs/heads/master/host/utils/uhd{-}usrp.rules }\KeywordTok{\&\&} \FunctionTok{cp}\NormalTok{ uhd{-}usrp.rules /etc/udev/rules.d/}
\end{Highlighting}
\end{Shaded}

\subsubsection{BladeRF (Nuand)}\label{bladerf-nuand}

\begin{itemize}
\tightlist
\item
  \textbf{libbladerf-dev} + build depuis les sources :
\end{itemize}

\begin{Shaded}
\begin{Highlighting}[]
\FunctionTok{git}\NormalTok{ clone https://github.com/nuand/BladeRF /opt/GSM/BladeRF}
\BuiltInTok{cd}\NormalTok{ /opt/GSM/BladeRF }\KeywordTok{\&\&} \FunctionTok{mkdir}\NormalTok{ build }\KeywordTok{\&\&} \BuiltInTok{cd}\NormalTok{ build}
\FunctionTok{cmake}\NormalTok{ .. }\KeywordTok{\&\&} \FunctionTok{make} \AttributeTok{{-}j}\VariableTok{$(}\FunctionTok{nproc}\VariableTok{)} \KeywordTok{\&\&} \FunctionTok{make}\NormalTok{ install }\KeywordTok{\&\&} \ExtensionTok{ldconfig}
\end{Highlighting}
\end{Shaded}

\begin{center}\rule{0.5\linewidth}{0.5pt}\end{center}

\subsection{Librairies Osmocom de
base}\label{librairies-osmocom-de-base}

Ces librairies sont la fondation de toute la stack GSM Osmocom.

\begin{itemize}
\tightlist
\item
  \textbf{libosmocore} : Primitives GSM, timers, gestion protocoles.
\item
  \textbf{libosmo-netif} : Interfaces réseau.
\item
  \textbf{libosmo-abis} : Support interface Abis (BTS ↔ BSC).
\end{itemize}

\textbf{Build de chaque librairie :}

\begin{Shaded}
\begin{Highlighting}[]
\FunctionTok{git}\NormalTok{ clone https://github.com/osmocom/}\OperatorTok{\textless{}}\NormalTok{repo}\OperatorTok{\textgreater{}}
\BuiltInTok{cd} \OperatorTok{\textless{}}\NormalTok{repo}\OperatorTok{\textgreater{}} \KeywordTok{\&\&} \ExtensionTok{autoreconf} \AttributeTok{{-}fi} \KeywordTok{\&\&} \ExtensionTok{./configure} \KeywordTok{\&\&} \FunctionTok{make} \AttributeTok{{-}j}\VariableTok{$(}\FunctionTok{nproc}\VariableTok{)} \KeywordTok{\&\&} \FunctionTok{make}\NormalTok{ install }\KeywordTok{\&\&} \ExtensionTok{ldconfig}
\end{Highlighting}
\end{Shaded}

\emph{(remplacer \texttt{\textless{}repo\textgreater{}} par le nom de la
librairie)}

\begin{center}\rule{0.5\linewidth}{0.5pt}\end{center}

\subsection{Stack cœur GSM/EDGE}\label{stack-cux153ur-gsmedge}

\pandocbounded{\includegraphics[keepaspectratio]{pictures/osmocom.png}}

\textbf{Osmocom fournit un module séparé pour chaque composant cœur de
réseau :}

\begin{itemize}
\tightlist
\item
  \textbf{osmo-hlr} : Home Location Register (base d'abonnés)
\item
  \textbf{osmo-msc} : Mobile Switching Center (appels, SMS, mobilité)
\item
  \textbf{osmo-bsc} : Base Station Controller (gestion BTS, handover)
\item
  \textbf{osmo-sgsn} : Serving GPRS Support Node (sessions GPRS/EDGE)
\item
  \textbf{osmo-ggsn} : Gateway GPRS Support Node (passerelle IP/data)
\item
  \textbf{osmo-mgw} : Media Gateway (flux voix/data)
\item
  \textbf{osmo-pcu} : Packet Control Unit (gestion trames GPRS/EDGE)
\item
  \textbf{libosmo-sigtran} : SIGTRAN/SS7 over IP
\item
  \textbf{libsmpp34} : Librairie SMPP (SMS)
\item
  \textbf{libgtpnl} : GTP (GPRS Tunneling Protocol pour data mobile)
\end{itemize}

\begin{center}\includegraphics[width=1.71in]{../../../../tmp/RtmpVUNkQ4/file28a55171e778d} \end{center}

\begin{center}\includegraphics[width=1.53in]{../../../../tmp/RtmpVUNkQ4/file28a555511195} \end{center}

\textbf{Exemple de build générique :} \textbf{Exemple de build générique
:}

\begin{Shaded}
\begin{Highlighting}[]
\FunctionTok{git}\NormalTok{ clone https://github.com/osmocom/}\OperatorTok{\textless{}}\NormalTok{module}\OperatorTok{\textgreater{}}
\BuiltInTok{cd} \OperatorTok{\textless{}}\NormalTok{module}\OperatorTok{\textgreater{}} \KeywordTok{\&\&} \ExtensionTok{autoreconf} \AttributeTok{{-}fi} \KeywordTok{\&\&} \ExtensionTok{./configure} \PreprocessorTok{[}\SpecialStringTok{OPTIONS}\PreprocessorTok{]} \KeywordTok{\&\&} \FunctionTok{make} \AttributeTok{{-}j}\VariableTok{$(}\FunctionTok{nproc}\VariableTok{)} \KeywordTok{\&\&} \FunctionTok{make}\NormalTok{ install }\KeywordTok{\&\&} \ExtensionTok{ldconfig}
\end{Highlighting}
\end{Shaded}

\begin{itemize}
\tightlist
\item
  Les options comme \texttt{-\/-enable-gtp-linux},
  \texttt{-\/-enable-iu}, \texttt{-\/-enable-smpp}, etc, dépendent du
  module.
\end{itemize}

\begin{center}\rule{0.5\linewidth}{0.5pt}\end{center}

\subsection{Couches radio (TRX/BTS)}\label{couches-radio-trxbts}

\begin{itemize}
\item
  \textbf{osmo-trx} : Transceiver SDR (modulation, A5/x, couche radio
  PHY)
\item
  \textbf{osmo-bts} : Interface logicielle pour le TRX (logiciel BTS)

  \begin{itemize}
  \tightlist
  \item
    Option : \texttt{-\/-enable-trx} pour activer le mode SDR.
  \end{itemize}
\end{itemize}

\textbf{Build typique :}

\begin{Shaded}
\begin{Highlighting}[]
\FunctionTok{git}\NormalTok{ clone https://github.com/osmocom/osmo{-}trx}
\BuiltInTok{cd}\NormalTok{ osmo{-}trx }\KeywordTok{\&\&} \ExtensionTok{autoreconf} \AttributeTok{{-}fi} \KeywordTok{\&\&} \ExtensionTok{./configure} \AttributeTok{{-}{-}with{-}uhd} \KeywordTok{\&\&} \FunctionTok{make} \AttributeTok{{-}j}\VariableTok{$(}\FunctionTok{nproc}\VariableTok{)} \KeywordTok{\&\&} \FunctionTok{make}\NormalTok{ install }\KeywordTok{\&\&} \ExtensionTok{ldconfig}

\FunctionTok{git}\NormalTok{ clone https://github.com/osmocom/osmo{-}bts}
\BuiltInTok{cd}\NormalTok{ osmo{-}bts }\KeywordTok{\&\&} \ExtensionTok{autoreconf} \AttributeTok{{-}fi} \KeywordTok{\&\&} \ExtensionTok{./configure} \AttributeTok{{-}{-}enable{-}trx} \KeywordTok{\&\&} \FunctionTok{make} \AttributeTok{{-}j}\VariableTok{$(}\FunctionTok{nproc}\VariableTok{)} \KeywordTok{\&\&} \FunctionTok{make}\NormalTok{ install }\KeywordTok{\&\&} \ExtensionTok{ldconfig}
\end{Highlighting}
\end{Shaded}

\begin{center}\rule{0.5\linewidth}{0.5pt}\end{center}

\subsection{Services \& Outils
additionnels}\label{services-outils-additionnels}

\begin{itemize}
\tightlist
\item
  \textbf{osmo-sip-connector} : passerelle GSM ↔ VoIP SIP (pour appels
  IP)
\item
  \textbf{Outils SIM} : \texttt{pcsc-tools}, \texttt{pcscd},
  \texttt{libpcsclite-dev}
\item
  \textbf{iptables}, \textbf{tmux}, \textbf{nano}, \textbf{python3-pip}
  : utilitaires système
\item
  \textbf{python3-requests} : scripts API/réseau
\end{itemize}

\textbf{Exemple d'installation :}

\begin{Shaded}
\begin{Highlighting}[]
\ExtensionTok{apt}\NormalTok{ install }\AttributeTok{{-}y}\NormalTok{ iptables tmux nano python3{-}pip python3{-}requests}
\ExtensionTok{pip3}\NormalTok{ install requests}
\end{Highlighting}
\end{Shaded}

\begin{center}\rule{0.5\linewidth}{0.5pt}\end{center}

\subsection{Script modulaire (bash)}\label{script-modulaire-bash}

\begin{Shaded}
\begin{Highlighting}[]
\CommentTok{\#!/bin/bash}
\CommentTok{\# Installeur modulaire EGPRS/Osmocom}
\CommentTok{\# Construit une stack GSM/EDGE complète avec SDR : USRP, BladeRF, Osmocom core et toutes dépendances.}
\CommentTok{\# Auteur : Bastien / bbaranoff (software{-}defined{-}radio.com)}
\CommentTok{\# Testé sur Ubuntu 22.04+}

\BuiltInTok{set} \AttributeTok{{-}e}
\VariableTok{N}\OperatorTok{=}\VariableTok{$(}\FunctionTok{nproc}\VariableTok{)} \KeywordTok{||} \VariableTok{N}\OperatorTok{=}\NormalTok{2}
\VariableTok{ROOT}\OperatorTok{=}\NormalTok{/opt/GSM}

\CommentTok{\# {-}{-}{-}{-}{-}{-}{-}{-}{-}{-}{-}{-}{-}{-}{-}{-}{-}{-}{-}{-}{-}{-}{-}{-}{-}{-}{-}{-}{-}{-}{-}{-}{-}{-}{-}{-}{-}{-}{-}{-}}
\CommentTok{\#  Dépendances système}
\CommentTok{\# {-}{-}{-}{-}{-}{-}{-}{-}{-}{-}{-}{-}{-}{-}{-}{-}{-}{-}{-}{-}{-}{-}{-}{-}{-}{-}{-}{-}{-}{-}{-}{-}{-}{-}{-}{-}{-}{-}{-}{-}}
\FunctionTok{install\_deps()} \KeywordTok{\{}
    \BuiltInTok{echo} \StringTok{"==\textgreater{} Installation des paquets système et dépendances de build..."}
    \ExtensionTok{apt}\NormalTok{ update}
    \ExtensionTok{apt}\NormalTok{ install }\AttributeTok{{-}y} \DataTypeTok{\textbackslash{}}
\NormalTok{        gcc g++ make cmake autoconf automake libtool pkg{-}config }\DataTypeTok{\textbackslash{}}
\NormalTok{        libedit{-}dev libfftw3{-}dev libusb{-}1.0{-}0{-}dev liburing{-}dev }\DataTypeTok{\textbackslash{}}
\NormalTok{        nano git wget curl tmux doxygen ethtool }\DataTypeTok{\textbackslash{}}
\NormalTok{        python3 python3{-}dev python3{-}pip python3{-}six python3{-}mako python3{-}lxml python3{-}numpy python3{-}pybind11 python3{-}requests python3{-}scipy python3{-}setuptools python3{-}ruamel.yaml }\DataTypeTok{\textbackslash{}}
\NormalTok{        libpcsclite{-}dev pcsc{-}tools pcscd cpufrequtils libsctp{-}dev libconfig++{-}dev libconfig{-}dev libmbedtls{-}dev }\DataTypeTok{\textbackslash{}}
\NormalTok{        libtalloc{-}dev libgnutls28{-}dev libmnl{-}dev libortp{-}dev libdbd{-}sqlite3 iproute2 iputils{-}ping }\DataTypeTok{\textbackslash{}}
\NormalTok{        libsqlite3{-}dev libdbi{-}dev libc{-}ares{-}dev libsofia{-}sip{-}ua{-}glib{-}dev telnet }\DataTypeTok{\textbackslash{}}
\NormalTok{        systemd systemd{-}sysv dbus dbus{-}user{-}session ccache libboost{-}all{-}dev libncurses5 libncurses5{-}dev }\DataTypeTok{\textbackslash{}}
\NormalTok{        libusb{-}1.0{-}0 libusb{-}dev iptables udev}
\KeywordTok{\}}

\CommentTok{\# {-}{-}{-}{-}{-}{-}{-}{-}{-}{-}{-}{-}{-}{-}{-}{-}{-}{-}{-}{-}{-}{-}{-}{-}{-}{-}{-}{-}{-}{-}{-}{-}{-}{-}{-}{-}{-}{-}{-}{-}}
\CommentTok{\#  Drivers SDR \& Firmware}
\CommentTok{\# {-}{-}{-}{-}{-}{-}{-}{-}{-}{-}{-}{-}{-}{-}{-}{-}{-}{-}{-}{-}{-}{-}{-}{-}{-}{-}{-}{-}{-}{-}{-}{-}{-}{-}{-}{-}{-}{-}{-}{-}}

\FunctionTok{install\_uhd()} \KeywordTok{\{}
    \BuiltInTok{echo} \StringTok{"==\textgreater{} Installation UHD/USRP et téléchargement images firmware..."}
    \ExtensionTok{apt}\NormalTok{ install }\AttributeTok{{-}y}\NormalTok{ libuhd{-}dev uhd{-}host}
    \ExtensionTok{python3}\NormalTok{ /usr/lib/uhd/utils/uhd\_images\_downloader.py}
    \CommentTok{\# Règles udev Ettus USRP}
    \FunctionTok{wget} \AttributeTok{{-}O}\NormalTok{ /etc/udev/rules.d/uhd{-}usrp.rules https://raw.githubusercontent.com/EttusResearch/uhd/refs/heads/master/host/utils/uhd{-}usrp.rules}
    \BuiltInTok{echo} \StringTok{"export UHD\_MODULE\_PATH=/usr/lib/uhd/modules"} \OperatorTok{\textgreater{}\textgreater{}}\NormalTok{ /root/.bashrc}
\KeywordTok{\}}

\FunctionTok{install\_bladerf()} \KeywordTok{\{}
    \BuiltInTok{echo} \StringTok{"==\textgreater{} Installation BladeRF et build depuis sources..."}
    \ExtensionTok{apt}\NormalTok{ install }\AttributeTok{{-}y}\NormalTok{ libbladerf{-}dev}
    \FunctionTok{git}\NormalTok{ clone https://github.com/nuand/BladeRF }\VariableTok{$ROOT}\NormalTok{/BladeRF }\KeywordTok{||} \FunctionTok{true}
    \BuiltInTok{cd} \VariableTok{$ROOT}\NormalTok{/BladeRF}
    \FunctionTok{mkdir} \AttributeTok{{-}p}\NormalTok{ build }\KeywordTok{\&\&} \BuiltInTok{cd}\NormalTok{ build}
    \FunctionTok{cmake}\NormalTok{ .. }\KeywordTok{\&\&} \FunctionTok{make} \AttributeTok{{-}j}\VariableTok{$N} \KeywordTok{\&\&} \FunctionTok{make}\NormalTok{ install }\KeywordTok{\&\&} \ExtensionTok{ldconfig}
\KeywordTok{\}}

\CommentTok{\# {-}{-}{-}{-}{-}{-}{-}{-}{-}{-}{-}{-}{-}{-}{-}{-}{-}{-}{-}{-}{-}{-}{-}{-}{-}{-}{-}{-}{-}{-}{-}{-}{-}{-}{-}{-}{-}{-}{-}{-}}
\CommentTok{\# Librairies Osmocom de base}
\CommentTok{\# {-}{-}{-}{-}{-}{-}{-}{-}{-}{-}{-}{-}{-}{-}{-}{-}{-}{-}{-}{-}{-}{-}{-}{-}{-}{-}{-}{-}{-}{-}{-}{-}{-}{-}{-}{-}{-}{-}{-}{-}}
\FunctionTok{build\_osmo\_lib()} \KeywordTok{\{}
    \BuiltInTok{echo} \StringTok{"==\textgreater{} Build des librairies Osmocom : libosmocore, libosmo{-}netif, libosmo{-}abis"}
    \ControlFlowTok{for}\NormalTok{ repo }\KeywordTok{in}\NormalTok{ libosmocore libosmo{-}netif libosmo{-}abis}\KeywordTok{;} \ControlFlowTok{do}
        \FunctionTok{git}\NormalTok{ clone https://github.com/osmocom/}\VariableTok{$repo} \VariableTok{$ROOT}\NormalTok{/}\VariableTok{$repo} \KeywordTok{||} \FunctionTok{true}
        \BuiltInTok{cd} \VariableTok{$ROOT}\NormalTok{/}\VariableTok{$repo}
        \ExtensionTok{autoreconf} \AttributeTok{{-}fi}
        \ExtensionTok{./configure}
        \FunctionTok{make} \AttributeTok{{-}j}\VariableTok{$N}
        \FunctionTok{make}\NormalTok{ install}
        \ExtensionTok{ldconfig}
    \ControlFlowTok{done}
\KeywordTok{\}}

\CommentTok{\# {-}{-}{-}{-}{-}{-}{-}{-}{-}{-}{-}{-}{-}{-}{-}{-}{-}{-}{-}{-}{-}{-}{-}{-}{-}{-}{-}{-}{-}{-}{-}{-}{-}{-}{-}{-}{-}{-}{-}{-}}
\CommentTok{\# Réseau cœur, radio, EDGE}
\CommentTok{\# {-}{-}{-}{-}{-}{-}{-}{-}{-}{-}{-}{-}{-}{-}{-}{-}{-}{-}{-}{-}{-}{-}{-}{-}{-}{-}{-}{-}{-}{-}{-}{-}{-}{-}{-}{-}{-}{-}{-}{-}}
\FunctionTok{build\_osmo\_core()} \KeywordTok{\{}
    \BuiltInTok{echo} \StringTok{"==\textgreater{} Build Osmocom réseau, radio et modules GPRS..."}
    
    \CommentTok{\# HLR (DB abonnés)}
    \FunctionTok{git}\NormalTok{ clone https://github.com/osmocom/osmo{-}hlr }\VariableTok{$ROOT}\NormalTok{/osmo{-}hlr }\KeywordTok{||} \FunctionTok{true}
    \BuiltInTok{cd} \VariableTok{$ROOT}\NormalTok{/osmo{-}hlr }\KeywordTok{\&\&} \ExtensionTok{autoreconf} \AttributeTok{{-}fi} \KeywordTok{\&\&} \ExtensionTok{./configure} \KeywordTok{\&\&} \FunctionTok{make} \AttributeTok{{-}j}\VariableTok{$N} \KeywordTok{\&\&} \FunctionTok{make}\NormalTok{ install }\KeywordTok{\&\&} \ExtensionTok{ldconfig}

    \CommentTok{\# Media Gateway}
    \FunctionTok{git}\NormalTok{ clone https://github.com/osmocom/osmo{-}mgw }\VariableTok{$ROOT}\NormalTok{/osmo{-}mgw }\KeywordTok{||} \FunctionTok{true}
    \BuiltInTok{cd} \VariableTok{$ROOT}\NormalTok{/osmo{-}mgw }\KeywordTok{\&\&} \ExtensionTok{autoreconf} \AttributeTok{{-}fi} \KeywordTok{\&\&} \ExtensionTok{./configure} \KeywordTok{\&\&} \FunctionTok{make} \AttributeTok{{-}j}\VariableTok{$N} \KeywordTok{\&\&} \FunctionTok{make}\NormalTok{ install }\KeywordTok{\&\&} \ExtensionTok{ldconfig}

    \CommentTok{\# SIGTRAN (SS7 sur IP)}
    \FunctionTok{git}\NormalTok{ clone https://gitea.osmocom.org/osmocom/libosmo{-}sigtran }\VariableTok{$ROOT}\NormalTok{/libosmo{-}sigtran }\KeywordTok{||} \FunctionTok{true}
    \BuiltInTok{cd} \VariableTok{$ROOT}\NormalTok{/libosmo{-}sigtran }\KeywordTok{\&\&} \ExtensionTok{autoreconf} \AttributeTok{{-}fi} \KeywordTok{\&\&} \ExtensionTok{./configure} \KeywordTok{\&\&} \FunctionTok{make} \AttributeTok{{-}j}\VariableTok{$N} \KeywordTok{\&\&} \FunctionTok{make}\NormalTok{ install }\KeywordTok{\&\&} \ExtensionTok{ldconfig}

    \CommentTok{\# libsmpp34 (SMS)}
    \FunctionTok{git}\NormalTok{ clone https://github.com/osmocom/libsmpp34 }\VariableTok{$ROOT}\NormalTok{/libsmpp34 }\KeywordTok{||} \FunctionTok{true}
    \BuiltInTok{cd} \VariableTok{$ROOT}\NormalTok{/libsmpp34 }\KeywordTok{\&\&} \ExtensionTok{autoreconf} \AttributeTok{{-}fi} \KeywordTok{\&\&} \ExtensionTok{./configure} \KeywordTok{\&\&} \FunctionTok{make} \AttributeTok{{-}j}\VariableTok{$N} \KeywordTok{\&\&} \FunctionTok{make}\NormalTok{ install }\KeywordTok{\&\&} \ExtensionTok{ldconfig}

    \CommentTok{\# libgtpnl (tunnels GTP)}
    \FunctionTok{git}\NormalTok{ clone https://github.com/osmocom/libgtpnl }\VariableTok{$ROOT}\NormalTok{/libgtpnl }\KeywordTok{||} \FunctionTok{true}
    \BuiltInTok{cd} \VariableTok{$ROOT}\NormalTok{/libgtpnl }\KeywordTok{\&\&} \ExtensionTok{autoreconf} \AttributeTok{{-}fi} \KeywordTok{\&\&} \ExtensionTok{./configure} \AttributeTok{{-}{-}enable{-}gtp{-}linux} \KeywordTok{\&\&} \FunctionTok{make} \AttributeTok{{-}j}\VariableTok{$N} \KeywordTok{\&\&} \FunctionTok{make}\NormalTok{ install }\KeywordTok{\&\&} \ExtensionTok{ldconfig}

    \CommentTok{\# Éléments cœur (ordre : GGSN, SGSN, MSC, BSC, PCU)}
    \FunctionTok{git}\NormalTok{ clone https://github.com/osmocom/osmo{-}ggsn }\VariableTok{$ROOT}\NormalTok{/osmo{-}ggsn }\KeywordTok{||} \FunctionTok{true}
    \BuiltInTok{cd} \VariableTok{$ROOT}\NormalTok{/osmo{-}ggsn }\KeywordTok{\&\&} \ExtensionTok{autoreconf} \AttributeTok{{-}fi} \KeywordTok{\&\&} \ExtensionTok{./configure} \AttributeTok{{-}{-}enable{-}gtp{-}linux} \KeywordTok{\&\&} \FunctionTok{make} \AttributeTok{{-}j}\VariableTok{$N} \KeywordTok{\&\&} \FunctionTok{make}\NormalTok{ install }\KeywordTok{\&\&} \ExtensionTok{ldconfig}

    \FunctionTok{git}\NormalTok{ clone https://github.com/osmocom/osmo{-}sgsn }\VariableTok{$ROOT}\NormalTok{/osmo{-}sgsn }\KeywordTok{||} \FunctionTok{true}
    \BuiltInTok{cd} \VariableTok{$ROOT}\NormalTok{/osmo{-}sgsn }\KeywordTok{\&\&} \ExtensionTok{autoreconf} \AttributeTok{{-}fi} \KeywordTok{\&\&} \ExtensionTok{./configure} \AttributeTok{{-}{-}enable{-}iu} \KeywordTok{\&\&} \FunctionTok{make} \AttributeTok{{-}j}\VariableTok{$N} \KeywordTok{\&\&} \FunctionTok{make}\NormalTok{ install }\KeywordTok{\&\&} \ExtensionTok{ldconfig}

    \FunctionTok{git}\NormalTok{ clone https://github.com/osmocom/osmo{-}msc }\VariableTok{$ROOT}\NormalTok{/osmo{-}msc }\KeywordTok{||} \FunctionTok{true}
    \BuiltInTok{cd} \VariableTok{$ROOT}\NormalTok{/osmo{-}msc }\KeywordTok{\&\&} \ExtensionTok{autoreconf} \AttributeTok{{-}fi} \KeywordTok{\&\&} \ExtensionTok{./configure} \AttributeTok{{-}{-}enable{-}smpp} \AttributeTok{{-}{-}enable{-}iu} \KeywordTok{\&\&} \FunctionTok{make} \AttributeTok{{-}j}\VariableTok{$N} \KeywordTok{\&\&} \FunctionTok{make}\NormalTok{ install }\KeywordTok{\&\&} \ExtensionTok{ldconfig}

    \FunctionTok{git}\NormalTok{ clone https://github.com/osmocom/osmo{-}bsc }\VariableTok{$ROOT}\NormalTok{/osmo{-}bsc }\KeywordTok{||} \FunctionTok{true}
    \BuiltInTok{cd} \VariableTok{$ROOT}\NormalTok{/osmo{-}bsc }\KeywordTok{\&\&} \ExtensionTok{autoreconf} \AttributeTok{{-}fi} \KeywordTok{\&\&} \ExtensionTok{./configure} \AttributeTok{{-}{-}enable{-}iu} \KeywordTok{\&\&} \FunctionTok{make} \AttributeTok{{-}j}\VariableTok{$N} \KeywordTok{\&\&} \FunctionTok{make}\NormalTok{ install }\KeywordTok{\&\&} \ExtensionTok{ldconfig}

    \FunctionTok{git}\NormalTok{ clone https://github.com/osmocom/osmo{-}pcu }\VariableTok{$ROOT}\NormalTok{/osmo{-}pcu }\KeywordTok{||} \FunctionTok{true}
    \BuiltInTok{cd} \VariableTok{$ROOT}\NormalTok{/osmo{-}pcu }\KeywordTok{\&\&} \ExtensionTok{autoreconf} \AttributeTok{{-}fi} \KeywordTok{\&\&} \ExtensionTok{./configure} \KeywordTok{\&\&} \FunctionTok{make} \AttributeTok{{-}j}\VariableTok{$N} \KeywordTok{\&\&} \FunctionTok{make}\NormalTok{ install }\KeywordTok{\&\&} \ExtensionTok{ldconfig}

    \CommentTok{\# Couches radio : TRX, BTS}
    \FunctionTok{git}\NormalTok{ clone https://github.com/osmocom/osmo{-}trx }\VariableTok{$ROOT}\NormalTok{/osmo{-}trx }\KeywordTok{||} \FunctionTok{true}
    \BuiltInTok{cd} \VariableTok{$ROOT}\NormalTok{/osmo{-}trx }\KeywordTok{\&\&} \ExtensionTok{autoreconf} \AttributeTok{{-}fi} \KeywordTok{\&\&} \ExtensionTok{./configure} \AttributeTok{{-}{-}with{-}uhd} \KeywordTok{\&\&} \FunctionTok{make} \AttributeTok{{-}j}\VariableTok{$N} \KeywordTok{\&\&} \FunctionTok{make}\NormalTok{ install }\KeywordTok{\&\&} \ExtensionTok{ldconfig}

    \FunctionTok{git}\NormalTok{ clone https://github.com/osmocom/osmo{-}bts }\VariableTok{$ROOT}\NormalTok{/osmo{-}bts }\KeywordTok{||} \FunctionTok{true}
    \BuiltInTok{cd} \VariableTok{$ROOT}\NormalTok{/osmo{-}bts }\KeywordTok{\&\&} \ExtensionTok{autoreconf} \AttributeTok{{-}fi} \KeywordTok{\&\&} \ExtensionTok{./configure} \AttributeTok{{-}{-}enable{-}trx} \KeywordTok{\&\&} \FunctionTok{make} \AttributeTok{{-}j}\VariableTok{$N} \KeywordTok{\&\&} \FunctionTok{make}\NormalTok{ install }\KeywordTok{\&\&} \ExtensionTok{ldconfig}

    \CommentTok{\# SIP connector (SIP/VoIP)}
    \FunctionTok{git}\NormalTok{ clone https://github.com/osmocom/osmo{-}sip{-}connector }\VariableTok{$ROOT}\NormalTok{/osmo{-}sip{-}connector }\KeywordTok{||} \FunctionTok{true}
    \BuiltInTok{cd} \VariableTok{$ROOT}\NormalTok{/osmo{-}sip{-}connector }\KeywordTok{\&\&} \ExtensionTok{autoreconf} \AttributeTok{{-}fi} \KeywordTok{\&\&} \ExtensionTok{./configure} \KeywordTok{\&\&} \FunctionTok{make} \AttributeTok{{-}j}\VariableTok{$N} \KeywordTok{\&\&} \FunctionTok{make}\NormalTok{ install }\KeywordTok{\&\&} \ExtensionTok{ldconfig}
\KeywordTok{\}}

\CommentTok{\# {-}{-}{-}{-}{-}{-}{-}{-}{-}{-}{-}{-}{-}{-}{-}{-}{-}{-}{-}{-}{-}{-}{-}{-}{-}{-}{-}{-}{-}{-}{-}{-}{-}{-}{-}{-}{-}{-}{-}{-}}
\CommentTok{\# Python \& utilitaires système}
\CommentTok{\# {-}{-}{-}{-}{-}{-}{-}{-}{-}{-}{-}{-}{-}{-}{-}{-}{-}{-}{-}{-}{-}{-}{-}{-}{-}{-}{-}{-}{-}{-}{-}{-}{-}{-}{-}{-}{-}{-}{-}{-}}
\FunctionTok{install\_python\_utils()} \KeywordTok{\{}
    \BuiltInTok{echo} \StringTok{"==\textgreater{} Installation dépendances pip et update images UHD..."}
    \ExtensionTok{pip3}\NormalTok{ install requests}
    \CommentTok{\# (Re)télécharge images firmware UHD si besoin}
    \ExtensionTok{python3}\NormalTok{ /usr/lib/uhd/utils/uhd\_images\_downloader.py }\KeywordTok{||} \FunctionTok{true}
\KeywordTok{\}}

\CommentTok{\# {-}{-}{-}{-}{-}{-}{-}{-}{-}{-}{-}{-}{-}{-}{-}{-}{-}{-}{-}{-}{-}{-}{-}{-}{-}{-}{-}{-}{-}{-}{-}{-}{-}{-}{-}{-}{-}{-}{-}{-}}
\CommentTok{\# Séquence principale}
\CommentTok{\# {-}{-}{-}{-}{-}{-}{-}{-}{-}{-}{-}{-}{-}{-}{-}{-}{-}{-}{-}{-}{-}{-}{-}{-}{-}{-}{-}{-}{-}{-}{-}{-}{-}{-}{-}{-}{-}{-}{-}{-}}
\FunctionTok{main()} \KeywordTok{\{}
    \ExtensionTok{install\_deps}
    \ExtensionTok{install\_uhd}
    \ExtensionTok{install\_bladerf}
    \ExtensionTok{build\_osmo\_lib}
    \ExtensionTok{build\_osmo\_core}
    \ExtensionTok{install\_python\_utils}
    \BuiltInTok{echo} \StringTok{"==\textgreater{} Tous les modules sont installés. Vous pouvez passer à la configuration et aux scripts de run."}
\KeywordTok{\}}

\ExtensionTok{main} \StringTok{"}\VariableTok{$@}\StringTok{"}
\end{Highlighting}
\end{Shaded}

\begin{center}\rule{0.5\linewidth}{0.5pt}\end{center}

\subsection{Résumé \& Bonnes
pratiques}\label{ruxe9sumuxe9-bonnes-pratiques}

\begin{itemize}
\tightlist
\item
  \textbf{Chaque module} est un bloc indépendant, à construire dans
  l'ordre (libs → cœur → radio → services)
\item
  Pour EGPRS, l'ordre \textbf{BTS/TRX/PCU/SGSN/GGSN} est crucial (le
  flux data traverse chaque élément)
\item
  \textbf{Séparation des modules} = maintenance et debug facilités
\item
  Toujours vérifier les dépendances système/matériel (drivers USRP,
  firmware BladeRF, règles udev\ldots)
\end{itemize}

\subsection{Installation d'une BTS OsmocomBB
Calypso}\label{installation-dune-bts-osmocombb-calypso}

\subsection{Matériel requis}\label{matuxe9riel-requis}

\begin{itemize}
\tightlist
\item
  1-2x Motorola C1xx (C115, C118, C123\ldots{} chipset Calypso)
\item
  2x Câbles USB-série jack 2.5\,mm (PL2303/FTDI)
\item
  Ubuntu 24.04 x86\_64 (VM ou natif)
\item
  Accès root
\end{itemize}

\begin{center}\rule{0.5\linewidth}{0.5pt}\end{center}

\subsection{Dépendances système}\label{duxe9pendances-systuxe8me-1}

\begin{Shaded}
\begin{Highlighting}[]
\FunctionTok{sudo}\NormalTok{ apt update}
\FunctionTok{sudo}\NormalTok{ apt install }\AttributeTok{{-}y} \DataTypeTok{\textbackslash{}}
\NormalTok{  build{-}essential git cmake pkg{-}config autoconf automake libtool shtool }\DataTypeTok{\textbackslash{}}
\NormalTok{  bison flex libtalloc{-}dev libpcsclite{-}dev zlib1g{-}dev libmpfr{-}dev libmpc{-}dev }\DataTypeTok{\textbackslash{}}
\NormalTok{  libgmp{-}dev libncurses5{-}dev libncursesw5{-}dev libsofia{-}sip{-}ua{-}dev libxml2{-}dev }\DataTypeTok{\textbackslash{}}
\NormalTok{  libfftw3{-}dev libgnutls28{-}dev libssl{-}dev python3 python3{-}pip }\DataTypeTok{\textbackslash{}}
\NormalTok{  alsa{-}oss lemon libtinfo{-}dev}
\end{Highlighting}
\end{Shaded}

\begin{quote}
\textbf{Note}\,: Ubuntu 24.04 a GCC\,13+ par défaut, mais il faut
parfois cross-compiler avec GCC 4.9 pour le firmware Calypso. S'il
existe des binaires arm-elf précompilés pour Ubuntu 22.04, ils
fonctionnent aussi sous 24.04. Pour l'ARM toolchain, voir plus bas
(utilise de préférence un container/VM si souci d'incompatibilité).
\end{quote}

\begin{center}\rule{0.5\linewidth}{0.5pt}\end{center}

\subsection{Installer la toolchain ARM-elf (cross-compiler
Calypso)}\label{installer-la-toolchain-arm-elf-cross-compiler-calypso}

\subsubsection{Option A\,: Téléchargement direct (prêt à
l'emploi)}\label{option-a-tuxe9luxe9chargement-direct-pruxeat-uxe0-lemploi}

\begin{Shaded}
\begin{Highlighting}[]
\CommentTok{\# Depuis Osmocom (préconisé : ne pas compiler vous{-}même sous Ubuntu récent)}
\BuiltInTok{cd}\NormalTok{ /opt}
\FunctionTok{sudo}\NormalTok{ wget https://downloads.osmocom.org/binaries/toolchain/arm{-}2021.03.tar.bz2}
\FunctionTok{sudo}\NormalTok{ tar }\AttributeTok{{-}xjf}\NormalTok{ arm{-}2021.03.tar.bz2}
\BuiltInTok{export} \VariableTok{PATH}\OperatorTok{=}\VariableTok{$PATH}\NormalTok{:/opt/arm{-}2021.03/bin}
\end{Highlighting}
\end{Shaded}

\begin{quote}
\textbf{Ajoute l'export PATH à ton .bashrc pour persistance}
\end{quote}

\subsubsection{Option B\,: Compilation
manuelle}\label{option-b-compilation-manuelle}

\begin{quote}
(non recommandé sur Ubuntu 24.04, c'est lent et casse souvent, préfère
la méthode ci-dessus)
\end{quote}

\begin{center}\rule{0.5\linewidth}{0.5pt}\end{center}

\subsection{Compiler et installer
OsmocomBB}\label{compiler-et-installer-osmocombb}

\begin{Shaded}
\begin{Highlighting}[]
\CommentTok{\# Cloner les sources}
\FunctionTok{git}\NormalTok{ clone https://github.com/osmocom/osmocom{-}bb.git}
\BuiltInTok{cd}\NormalTok{ osmocom{-}bb}

\CommentTok{\# Préparer le build}
\BuiltInTok{cd}\NormalTok{ src}
\ExtensionTok{autoreconf} \AttributeTok{{-}i}

\CommentTok{\# Compiler le firmware (spécifie la toolchain si besoin)}
\BuiltInTok{export} \VariableTok{PATH}\OperatorTok{=}\VariableTok{$PATH}\NormalTok{:/opt/arm{-}2021.03/bin}
\ExtensionTok{./configure}
\FunctionTok{make} \AttributeTok{{-}j}\VariableTok{$(}\FunctionTok{nproc}\VariableTok{)}
\FunctionTok{sudo}\NormalTok{ make install}
\FunctionTok{sudo}\NormalTok{ ldconfig}
\end{Highlighting}
\end{Shaded}

\begin{center}\rule{0.5\linewidth}{0.5pt}\end{center}

\subsection{Flasher le firmware sur les
téléphones}\label{flasher-le-firmware-sur-les-tuxe9luxe9phones}

\textbf{Connecte les deux téléphones via USB-PL2303, batterie retirée.}

\begin{Shaded}
\begin{Highlighting}[]
\BuiltInTok{cd}\NormalTok{ src/host/osmocon}
\FunctionTok{sudo}\NormalTok{ ./osmocon }\AttributeTok{{-}m}\NormalTok{ c123xor }\AttributeTok{{-}p}\NormalTok{ /dev/ttyUSB0 }\AttributeTok{{-}c}\NormalTok{ ../../target/firmware/board/compal\_e88/rssi.highram.bin}
\CommentTok{\# (Shell \#1 : RSSI scan, contrôle{-}C pour arrêter)}
\end{Highlighting}
\end{Shaded}

\textbf{Trouve un ARFCN actif (fort signal GSM, ex.~: 900/1800 MHz).}

\begin{Shaded}
\begin{Highlighting}[]
\CommentTok{\# Pour charger le firmware TRX (BTS)}
\FunctionTok{sudo}\NormalTok{ ./osmocon }\AttributeTok{{-}m}\NormalTok{ c123xor }\AttributeTok{{-}p}\NormalTok{ /dev/ttyUSB0 }\AttributeTok{{-}s}\NormalTok{ /tmp/osmocom\_l2 }\DataTypeTok{\textbackslash{}}
  \AttributeTok{{-}c}\NormalTok{ ../../target/firmware/board/compal\_e88/trx.highram.bin }\AttributeTok{{-}r}\NormalTok{ 99}
\FunctionTok{sudo}\NormalTok{ ./osmocon }\AttributeTok{{-}m}\NormalTok{ c123xor }\AttributeTok{{-}p}\NormalTok{ /dev/ttyUSB1 }\AttributeTok{{-}s}\NormalTok{ /tmp/osmocom\_l2.2 }\DataTypeTok{\textbackslash{}}
  \AttributeTok{{-}c}\NormalTok{ ../../target/firmware/board/compal\_e88/trx.highram.bin }\AttributeTok{{-}r}\NormalTok{ 99}
\end{Highlighting}
\end{Shaded}

\begin{center}\rule{0.5\linewidth}{0.5pt}\end{center}

\subsection{Lancer le transceiver et la stack
BTS}\label{lancer-le-transceiver-et-la-stack-bts}

Dans d'autres shells\,:

\begin{Shaded}
\begin{Highlighting}[]
\CommentTok{\# Shell \#3 : transceiver Calypso (remplace ARFCN trouvé)}
\BuiltInTok{cd}\NormalTok{ src/host/layer23/src/transceiver/}
\FunctionTok{sudo}\NormalTok{ ./transceiver }\AttributeTok{{-}a} \PreprocessorTok{[}\SpecialStringTok{ARFCN}\PreprocessorTok{]} \AttributeTok{{-}2} \AttributeTok{{-}r}\NormalTok{ 99}

\CommentTok{\# Shell \#4 : osmo{-}nitb (NITB = core minimal)}
\ExtensionTok{osmo{-}nitb} \AttributeTok{{-}c}\NormalTok{ \textasciitilde{}/.osmocom/open{-}bsc.cfg }\AttributeTok{{-}l}\NormalTok{ \textasciitilde{}/.osmocom/hlr.sqlite3 }\AttributeTok{{-}P} \AttributeTok{{-}m} \AttributeTok{{-}C} \AttributeTok{{-}{-}debug}\OperatorTok{=}\NormalTok{DRLL:DCC:DMM:DRR:DRSL:DNM}

\CommentTok{\# Shell \#5 : osmobts{-}trx}
\ExtensionTok{osmobts{-}trx} \AttributeTok{{-}c}\NormalTok{ \textasciitilde{}/.osmocom/osmo{-}bts.cfg }\AttributeTok{{-}r}\NormalTok{ 99}
\end{Highlighting}
\end{Shaded}

\begin{center}\rule{0.5\linewidth}{0.5pt}\end{center}

\subsection{(Optionnel) Installer Asterisk + LCR pour appels
SIP}\label{optionnel-installer-asterisk-lcr-pour-appels-sip}

\begin{Shaded}
\begin{Highlighting}[]
\FunctionTok{sudo}\NormalTok{ apt install }\AttributeTok{{-}y}\NormalTok{ asterisk asterisk{-}dev}
\FunctionTok{git}\NormalTok{ clone https://github.com/fairwaves/lcr}
\BuiltInTok{cd}\NormalTok{ lcr}
\ExtensionTok{autoreconf} \AttributeTok{{-}i}
\ExtensionTok{./configure} \AttributeTok{{-}{-}with{-}sip} \AttributeTok{{-}{-}with{-}gsm{-}bs} \AttributeTok{{-}{-}with{-}gsm{-}ms} \AttributeTok{{-}{-}with{-}asterisk}
\FunctionTok{make}
\FunctionTok{sudo}\NormalTok{ make install}
\FunctionTok{sudo}\NormalTok{ cp chan\_lcr.so /usr/lib/asterisk/modules/}
\end{Highlighting}
\end{Shaded}

Configurer \texttt{/etc/asterisk/sip.conf} avec tes identifiants SIP.

\begin{center}\rule{0.5\linewidth}{0.5pt}\end{center}

\subsection{Vérifier le setup et
logs}\label{vuxe9rifier-le-setup-et-logs}

\begin{itemize}
\tightlist
\item
  Vérifie les logs dans chaque shell\,: recherche ``registered'',
  ``attach'', etc.
\item
  Lancer l'attachement avec un mobile (SIM Osmocom ou Free avec accès
  PLMN custom).
\end{itemize}

\begin{center}\rule{0.5\linewidth}{0.5pt}\end{center}

\section{Installation de srsRAN 4G (via PPA ou
sources)}\label{installation-de-srsran-4g-via-ppa-ou-sources}

\pandocbounded{\includegraphics[keepaspectratio]{pictures/srsran.png}}

\subsection{Installation rapide via PPA
(Ubuntu)}\label{installation-rapide-via-ppa-ubuntu}

\begin{Shaded}
\begin{Highlighting}[]
\FunctionTok{sudo}\NormalTok{ apt update}
\FunctionTok{sudo}\NormalTok{ apt install }\AttributeTok{{-}y}\NormalTok{ software{-}properties{-}common}
\FunctionTok{sudo}\NormalTok{ add{-}apt{-}repository ppa:srsran/ppa}
\FunctionTok{sudo}\NormalTok{ apt update}
\FunctionTok{sudo}\NormalTok{ apt install }\AttributeTok{{-}y}\NormalTok{ srsran}
\end{Highlighting}
\end{Shaded}

\textbf{Vérifie l'install :}

\begin{Shaded}
\begin{Highlighting}[]
\ExtensionTok{srsenb} \AttributeTok{{-}{-}version}
\end{Highlighting}
\end{Shaded}

\begin{center}\rule{0.5\linewidth}{0.5pt}\end{center}

\subsection{Compilation depuis les
sources}\label{compilation-depuis-les-sources}

\begin{Shaded}
\begin{Highlighting}[]
\FunctionTok{sudo}\NormalTok{ apt update}
\FunctionTok{sudo}\NormalTok{ apt install }\AttributeTok{{-}y}\NormalTok{ git cmake build{-}essential libfftw3{-}dev libmbedtls{-}dev }\DataTypeTok{\textbackslash{}}
\NormalTok{  libboost{-}program{-}options{-}dev libconfig++{-}dev libsctp{-}dev }\DataTypeTok{\textbackslash{}}
\NormalTok{  libuhd{-}dev libpcsclite{-}dev pkg{-}config python3 python3{-}pip}

\FunctionTok{git}\NormalTok{ clone https://github.com/srsran/srsRAN\_Project.git}
\BuiltInTok{cd}\NormalTok{ srsRAN\_Project}
\FunctionTok{mkdir}\NormalTok{ build }\KeywordTok{\&\&} \BuiltInTok{cd}\NormalTok{ build}
\FunctionTok{cmake}\NormalTok{ ../}
\FunctionTok{make} \AttributeTok{{-}j}\VariableTok{$(}\FunctionTok{nproc}\VariableTok{)}
\FunctionTok{sudo}\NormalTok{ make install}
\end{Highlighting}
\end{Shaded}

\begin{center}\rule{0.5\linewidth}{0.5pt}\end{center}

\subsection{Générer/Configurer une SIM (IMSI, Ki,
OPC)}\label{guxe9nuxe9rerconfigurer-une-sim-imsi-ki-opc}

Pour connecter un téléphone réel ou un UE logiciel, il faut une SIM dont
les paramètres sont connus du core \textbf{srsEPC}.

\subsection{a. Générer un profil SIM
standard}\label{a.-guxe9nuxe9rer-un-profil-sim-standard}

\begin{itemize}
\tightlist
\item
  \textbf{IMSI} : identifiant unique (ex: \texttt{001010123456789})
\item
  \textbf{Ki} : clé de 16 octets (ex:
  \texttt{00112233445566778899aabbccddeeff})
\item
  \textbf{OPC} : clé d'opérateur de 16 octets (ex:
  \texttt{0102030405060708090a0b0c0d0e0f10})
\item
  \textbf{MSISDN} (optionnel) : numéro de tel virtuel
\end{itemize}

Tu peux utiliser un générateur en ligne
(\href{https://free5gc.org/webconsole/sim.html}{Free5GC SIM Generator}),
ou générer aléatoirement\,:

\begin{Shaded}
\begin{Highlighting}[]
\CommentTok{\# Générer un IMSI au hasard}
\VariableTok{IMSI}\OperatorTok{=}\StringTok{"001010}\VariableTok{$(}\FunctionTok{shuf} \AttributeTok{{-}i}\NormalTok{ 100000000{-}999999999 }\AttributeTok{{-}n}\NormalTok{ 1}\VariableTok{)}\StringTok{"}
\CommentTok{\# Générer une clé Ki et OPC}
\VariableTok{KI}\OperatorTok{=}\VariableTok{$(}\ExtensionTok{openssl}\NormalTok{ rand }\AttributeTok{{-}hex}\NormalTok{ 16}\VariableTok{)}
\VariableTok{OPC}\OperatorTok{=}\VariableTok{$(}\ExtensionTok{openssl}\NormalTok{ rand }\AttributeTok{{-}hex}\NormalTok{ 16}\VariableTok{)}
\BuiltInTok{echo} \StringTok{"IMSI: }\VariableTok{$IMSI}\StringTok{"}
\BuiltInTok{echo} \StringTok{"KI:   }\VariableTok{$KI}\StringTok{"}
\BuiltInTok{echo} \StringTok{"OPC:  }\VariableTok{$OPC}\StringTok{"}
\end{Highlighting}
\end{Shaded}

\subsection{b. Ajouter l'abonné dans
srsEPC}\label{b.-ajouter-labonnuxe9-dans-srsepc}

Ouvre le fichier de conf des abonnés, typiquement :

\begin{Shaded}
\begin{Highlighting}[]
\FunctionTok{nano}\NormalTok{ /etc/srsran/epc\_user\_db.csv}
\end{Highlighting}
\end{Shaded}

Ou dans ton dossier local \texttt{srsRAN\_Project/config/user\_db.csv}

\textbf{Exemple d'entrée :}

\begin{verbatim}
IMSI,  Ki,                       OPC,                      MSISDN,     AMF,   SQN
001010123456789,00112233445566778899aabbccddeeff,0102030405060708090a0b0c0d0e0f10,1234567890,8000,000000
\end{verbatim}

\textbf{Champ obligatoire}\,: IMSI, Ki, OPC (les autres = valeur par
défaut si besoin).

\begin{center}\rule{0.5\linewidth}{0.5pt}\end{center}

\subsection{Programmer une carte SIM réelle
(optionnel)}\label{programmer-une-carte-sim-ruxe9elle-optionnel}

Si tu utilises une carte programmable (Sysmocom, MagicSIM, etc.)\,:

\begin{itemize}
\tightlist
\item
  Utilise \href{https://github.com/osmocom/pysim}{pySIM} (Python)
\item
  Lance\,:
\end{itemize}

\begin{Shaded}
\begin{Highlighting}[]
\ExtensionTok{pysim{-}prog.py} \AttributeTok{{-}{-}ki} \VariableTok{$KI} \AttributeTok{{-}{-}opc} \VariableTok{$OPC} \AttributeTok{{-}{-}imsi} \VariableTok{$IMSI}
\end{Highlighting}
\end{Shaded}

(Adapté selon ton lecteur et modèle de carte)

\begin{center}\rule{0.5\linewidth}{0.5pt}\end{center}

\subsection{Lancer srsENB (eNodeB) et srsEPC
(core)}\label{lancer-srsenb-enodeb-et-srsepc-core}

Dans deux terminaux\,:

\begin{Shaded}
\begin{Highlighting}[]
\CommentTok{\# Terminal 1 (eNodeB)}
\ExtensionTok{srsenb}\NormalTok{ enb.conf}

\CommentTok{\# Terminal 2 (EPC)}
\ExtensionTok{srsepc}\NormalTok{ epc.conf}
\end{Highlighting}
\end{Shaded}

\begin{center}\rule{0.5\linewidth}{0.5pt}\end{center}

\subsection{Tester l'attachement UE}\label{tester-lattachement-ue}

\begin{itemize}
\tightlist
\item
  Allume ton téléphone (SIM insérée, mode réseau ``4G/auto'')
\item
  Vérifie la détection dans les logs
\item
  Si besoin, ajuste le PLMN dans \texttt{enb.conf} et \texttt{epc.conf}
  pour matcher l'IMSI de la SIM (ex: MCC/MNC = 001/01)
\end{itemize}

\begin{center}\rule{0.5\linewidth}{0.5pt}\end{center}

\subsubsection{Docs utiles}\label{docs-utiles}

\begin{itemize}
\tightlist
\item
  \url{https://docs.srsran.com/projects/4g/en/latest/index.html}
\item
  \url{https://osmocom.org/projects/pysim/wiki}
\end{itemize}

\section{L'évolution 5G : rupture technologique, nouveaux
usages}\label{luxe9volution-5g-rupture-technologique-nouveaux-usages}

\subsection{Objectifs et rupture de la
5G}\label{objectifs-et-rupture-de-la-5g}

\begin{itemize}
\tightlist
\item
  \textbf{Débit x10 à x100 par rapport à la 4G} (jusqu'à 10 Gbps en
  pointe, 1 Gbps courant)
\item
  \textbf{Ultra faible latence} (\textless\,1 ms possible)
\item
  \textbf{Connectivité massive} : jusqu'à 1 million d'objets/km² (IoT,
  capteurs, véhicules)
\item
  \textbf{Fiabilité et disponibilité quasi temps réel} (industrie,
  santé, véhicules autonomes)
\item
  \textbf{Flexibilité réseau} : découpage ``network slicing'', chaque
  usage a son ``mini-réseau'' dédié
\end{itemize}

\begin{center}\rule{0.5\linewidth}{0.5pt}\end{center}

\subsection{Ruptures techniques
majeures}\label{ruptures-techniques-majeures}

\begin{itemize}
\tightlist
\item
  \textbf{Nouvelles bandes de fréquences} : ajout du spectre
  millimétrique (\textgreater24 GHz), plus de capacité mais portée plus
  courte
\item
  \textbf{Massive MIMO} : dizaines/centaines d'antennes pour booster
  débit et couverture
\item
  \textbf{Beamforming} : orientation dynamique du signal vers chaque
  utilisateur
\item
  \textbf{Core 5G tout IP/cloud natif} : virtualisation (NFV), fonctions
  réseau en software (SDN)
\item
  \textbf{Architecture Standalone (SA)} : cœur 5G indépendant, vs
  Non-Standalone (NSA) où la 5G s'appuie sur la 4G
\end{itemize}

\begin{center}\rule{0.5\linewidth}{0.5pt}\end{center}

\subsection{Nouveaux usages rendus
possibles}\label{nouveaux-usages-rendus-possibles}

\begin{itemize}
\tightlist
\item
  \textbf{Internet mobile ultra haut débit} (streaming 4K/8K, réalité
  augmentée/virtuelle, cloud gaming)
\item
  \textbf{Industrie 4.0} : automatisation, robotique, maintenance
  prédictive, réseaux privés d'usine
\item
  \textbf{Véhicules connectés et autonomes} : communication V2X temps
  réel
\item
  \textbf{Santé} : télémédecine en temps réel, chirurgie à distance
\item
  \textbf{Smart Cities et IoT massif} : gestion d'énergie, surveillance,
  transports, capteurs urbains
\end{itemize}

\begin{center}\rule{0.5\linewidth}{0.5pt}\end{center}

\subsection{Limites, critiques et
défis}\label{limites-critiques-et-duxe9fis}

\begin{itemize}
\tightlist
\item
  \textbf{Couverture inégale} (surtout en mmWave, portée faible, besoin
  de micro-cellules)
\item
  \textbf{Consommation énergétique} (optimisée par le slicing mais
  déploiement massif d'antennes)
\item
  \textbf{Nouveaux risques de sécurité} (surface d'attaque plus large,
  virtualisation du cœur réseau)
\item
  \textbf{Investissements lourds} (infrastructures, renouvellement
  terminaux)
\end{itemize}

\begin{center}\rule{0.5\linewidth}{0.5pt}\end{center}

\subsection{Tableau récapitulatif}\label{tableau-ruxe9capitulatif}

\begin{longtable}[]{@{}
  >{\raggedright\arraybackslash}p{(\linewidth - 10\tabcolsep) * \real{0.1020}}
  >{\raggedright\arraybackslash}p{(\linewidth - 10\tabcolsep) * \real{0.1939}}
  >{\raggedright\arraybackslash}p{(\linewidth - 10\tabcolsep) * \real{0.0918}}
  >{\raggedright\arraybackslash}p{(\linewidth - 10\tabcolsep) * \real{0.1429}}
  >{\raggedright\arraybackslash}p{(\linewidth - 10\tabcolsep) * \real{0.1429}}
  >{\raggedright\arraybackslash}p{(\linewidth - 10\tabcolsep) * \real{0.3265}}@{}}
\toprule\noalign{}
\begin{minipage}[b]{\linewidth}\raggedright
Génération
\end{minipage} & \begin{minipage}[b]{\linewidth}\raggedright
Débit max
\end{minipage} & \begin{minipage}[b]{\linewidth}\raggedright
Latence
\end{minipage} & \begin{minipage}[b]{\linewidth}\raggedright
Connexions/km²
\end{minipage} & \begin{minipage}[b]{\linewidth}\raggedright
Architecture
\end{minipage} & \begin{minipage}[b]{\linewidth}\raggedright
Usages clés
\end{minipage} \\
\midrule\noalign{}
\endhead
\bottomrule\noalign{}
\endlastfoot
4G & \textasciitilde100 Mbps -- 1 Gbps & \textasciitilde50 ms &
\textasciitilde100\,000 & Tout-IP & Vidéo, web, app mobile \\
\textbf{5G} & 1--10 Gbps & \textless\,1--10 ms & 1\,000\,000+ & Cloud,
slicing & IoT massif, VR/AR, Industrie 4.0 \\
\end{longtable}

\begin{center}\rule{0.5\linewidth}{0.5pt}\end{center}

\subsection{À retenir}\label{uxe0-retenir}

\begin{quote}
\textbf{La 5G n'est pas juste ``plus rapide''}\,: elle transforme
l'infrastructure, permet l'émergence de nouveaux usages critiques
(industrie, santé, véhicules), et prépare le terrain pour un monde
ultra-connecté, piloté en temps réel.
\end{quote}

\section{5G SA labo\,: srsRAN Project + Open5GS (Ubuntu
2024)}\label{g-sa-labo-srsran-project-open5gs-ubuntu-2024}

\subsubsection{Architecture 5G SA (Open5GS +
srsRAN)}\label{architecture-5g-sa-open5gs-srsran}

\subsection{Prérequis}\label{pruxe9requis-1}

\begin{itemize}
\tightlist
\item
  \textbf{OS\,:} Ubuntu 22.04 ou 24.04 (x86\_64)
\item
  \textbf{Matériel\,:} SDR (USRP B200/B210/X300/PlutoSDR/LimeSDR\ldots)
  \emph{OU} mode simulation (aucun matériel SDR requis)
\item
  \textbf{Packages de base\,:} git, cmake, build-essential, etc.
\end{itemize}

\begin{center}\rule{0.5\linewidth}{0.5pt}\end{center}

\subsection{Installer Open5GS (Core
5G)}\label{installer-open5gs-core-5g}

\pandocbounded{\includegraphics[keepaspectratio]{pictures/open5gs.png}}

\subsubsection{a) Dépendances}\label{a-duxe9pendances}

\begin{Shaded}
\begin{Highlighting}[]
\FunctionTok{sudo}\NormalTok{ apt update}
\FunctionTok{sudo}\NormalTok{ apt install }\AttributeTok{{-}y}\NormalTok{ git build{-}essential meson ninja{-}build python3 python3{-}pip python3{-}setuptools }\DataTypeTok{\textbackslash{}}
\NormalTok{  flex bison libsctp{-}dev libssl{-}dev libnghttp2{-}dev libmicrohttpd{-}dev }\DataTypeTok{\textbackslash{}}
\NormalTok{  libyaml{-}dev libmariadb{-}dev libmongoc{-}dev libbson{-}dev mongodb}
\end{Highlighting}
\end{Shaded}

\subsubsection{b) Cloner et compiler
Open5GS}\label{b-cloner-et-compiler-open5gs}

\begin{Shaded}
\begin{Highlighting}[]
\FunctionTok{git}\NormalTok{ clone https://github.com/open5gs/open5gs}
\BuiltInTok{cd}\NormalTok{ open5gs}
\ExtensionTok{meson}\NormalTok{ build }\AttributeTok{{-}{-}prefix}\OperatorTok{=}\KeywordTok{\textasciigrave{}}\BuiltInTok{pwd}\KeywordTok{\textasciigrave{}}\NormalTok{/install}
\ExtensionTok{ninja} \AttributeTok{{-}C}\NormalTok{ build}
\ExtensionTok{ninja} \AttributeTok{{-}C}\NormalTok{ build install}
\end{Highlighting}
\end{Shaded}

\emph{(Binaires dans \texttt{open5gs/install/bin/})}

\begin{center}\rule{0.5\linewidth}{0.5pt}\end{center}

\subsection{Installer srsRAN Project
(gNB/UE)}\label{installer-srsran-project-gnbue}

\subsubsection{a) Dépendances}\label{a-duxe9pendances-1}

\begin{Shaded}
\begin{Highlighting}[]
\FunctionTok{sudo}\NormalTok{ apt install }\AttributeTok{{-}y}\NormalTok{ cmake g++ libsctp{-}dev libfftw3{-}dev libmbedtls{-}dev }\DataTypeTok{\textbackslash{}}
\NormalTok{  libboost{-}program{-}options{-}dev libconfig++{-}dev pkg{-}config python3 python3{-}pip }\DataTypeTok{\textbackslash{}}
\NormalTok{  libzmq3{-}dev libpcsclite{-}dev}
\end{Highlighting}
\end{Shaded}

\subsubsection{b) Cloner et compiler
srsRAN}\label{b-cloner-et-compiler-srsran}

\begin{Shaded}
\begin{Highlighting}[]
\FunctionTok{git}\NormalTok{ clone https://github.com/srsran/srsRAN\_Project.git}
\BuiltInTok{cd}\NormalTok{ srsRAN\_Project}
\FunctionTok{mkdir}\NormalTok{ build }\KeywordTok{\&\&} \BuiltInTok{cd}\NormalTok{ build}
\FunctionTok{cmake}\NormalTok{ ../}
\FunctionTok{make} \AttributeTok{{-}j}\VariableTok{$(}\FunctionTok{nproc}\VariableTok{)}
\FunctionTok{sudo}\NormalTok{ make install}
\end{Highlighting}
\end{Shaded}

\emph{(Binaires dans \texttt{/usr/local/bin/})}

\begin{center}\rule{0.5\linewidth}{0.5pt}\end{center}

\subsection{Configurer Open5GS (Core)}\label{configurer-open5gs-core}

\begin{itemize}
\item
  Tous les fichiers\,: \texttt{open5gs/install/etc/open5gs/}
\item
  Fichiers clés\,:

  \begin{itemize}
  \tightlist
  \item
    \texttt{amf.yaml} (AMF/core)
  \item
    \texttt{smf.yaml} (SMF/session)
  \item
    \texttt{upf.yaml} (UPF/user-plane)
  \item
    \texttt{udm.yaml}, \texttt{ausf.yaml}, etc.
  \end{itemize}
\end{itemize}

\textbf{Edits minimaux\,:}

\begin{itemize}
\tightlist
\item
  PLMN\,: \texttt{mcc:\ 001}, \texttt{mnc:\ 01} (ou tes propres valeurs)
\item
  Dans \texttt{amf.yaml}\,: l'adresse \texttt{ngap} doit matcher celle
  de ton gNB (ex\,: \texttt{127.0.0.1})
\end{itemize}

\subsubsection{Ajouter un abonné test
(SIM)}\label{ajouter-un-abonnuxe9-test-sim}

\begin{itemize}
\item
  Ouvre \texttt{subscribers.yaml} ou passe par le
  \href{https://github.com/open5gs/open5gs/wiki/WebUI}{WebUI Open5GS}
\item
  Exemple d'entrée :

\begin{Shaded}
\begin{Highlighting}[]
\KeywordTok{{-}}\AttributeTok{ }\FunctionTok{imsi}\KeywordTok{:}\AttributeTok{ }\StringTok{\textquotesingle{}001010123456789\textquotesingle{}}
\AttributeTok{  }\FunctionTok{key}\KeywordTok{:}\AttributeTok{ }\StringTok{\textquotesingle{}465B5CE8B199B49FAA5F0A2EE238A6BC\textquotesingle{}}
\AttributeTok{  }\FunctionTok{opc}\KeywordTok{:}\AttributeTok{ }\StringTok{\textquotesingle{}E8ED289DEBA952E4283B54E88E6183CA\textquotesingle{}}
\AttributeTok{  }\FunctionTok{amf}\KeywordTok{:}\AttributeTok{ }\StringTok{\textquotesingle{}8000\textquotesingle{}}
\AttributeTok{  }\FunctionTok{sqn}\KeywordTok{:}\AttributeTok{ }\StringTok{\textquotesingle{}000000000000\textquotesingle{}}
\end{Highlighting}
\end{Shaded}
\item
  À reporter dans la SIM (physique ou config UE logicielle)
\end{itemize}

\begin{center}\rule{0.5\linewidth}{0.5pt}\end{center}

\subsection{Configurer srsRAN gNB}\label{configurer-srsran-gnb}

\begin{itemize}
\tightlist
\item
  Fichier\,: \texttt{srsran\_project/config/gnb.yaml} \emph{(ou
  \texttt{gnb.conf} selon la version)}
\end{itemize}

Exemple minimal\,:

\begin{Shaded}
\begin{Highlighting}[]
\FunctionTok{plmn}\KeywordTok{:}
\AttributeTok{  }\FunctionTok{mcc}\KeywordTok{:}\AttributeTok{ }\DecValTok{001}
\AttributeTok{  }\FunctionTok{mnc}\KeywordTok{:}\AttributeTok{ }\DecValTok{01}
\FunctionTok{amf}\KeywordTok{:}
\AttributeTok{  }\FunctionTok{addr}\KeywordTok{:}\AttributeTok{ }\FloatTok{127.0.0.1}\CommentTok{   \# IP du core Open5GS}
\AttributeTok{  }\FunctionTok{port}\KeywordTok{:}\AttributeTok{ }\DecValTok{38412}
\FunctionTok{slices}\KeywordTok{:}
\AttributeTok{  }\KeywordTok{{-}}\AttributeTok{ }\FunctionTok{sst}\KeywordTok{:}\AttributeTok{ }\DecValTok{1}
\AttributeTok{    }\FunctionTok{sd}\KeywordTok{:}\AttributeTok{ }\DecValTok{0x000001}
\end{Highlighting}
\end{Shaded}

\begin{itemize}
\tightlist
\item
  Spécifie \texttt{rf.device\_name} pour ton SDR (si RF)
\item
  Réseau/interface\,: ex\,: \texttt{ngap.interface\_name:\ lo} pour
  tests locaux
\end{itemize}

\begin{center}\rule{0.5\linewidth}{0.5pt}\end{center}

\subsection{Configurer srsRAN UE (si test en pur
logiciel)}\label{configurer-srsran-ue-si-test-en-pur-logiciel}

\begin{itemize}
\tightlist
\item
  Fichier\,: \texttt{srsran\_project/config/ue.yaml} \emph{(ou
  \texttt{ue.conf})}
\item
  Paramètres\,: PLMN, key, opc, amf, imei identiques à l'abonné Open5GS.
\item
  Mode RF ou simulation selon besoin.
\end{itemize}

\begin{center}\rule{0.5\linewidth}{0.5pt}\end{center}

\subsection{Lancer le Core Open5GS}\label{lancer-le-core-open5gs}

Ouvre \textbf{plusieurs terminaux} (ou \texttt{tmux})\,:

\begin{Shaded}
\begin{Highlighting}[]
\BuiltInTok{cd}\NormalTok{ open5gs/install/bin/}
\FunctionTok{sudo}\NormalTok{ ./open5gs{-}mmed}
\FunctionTok{sudo}\NormalTok{ ./open5gs{-}sgwcd}
\FunctionTok{sudo}\NormalTok{ ./open5gs{-}smfd}
\FunctionTok{sudo}\NormalTok{ ./open5gs{-}amfd}
\FunctionTok{sudo}\NormalTok{ ./open5gs{-}upfd}
\FunctionTok{sudo}\NormalTok{ ./open5gs{-}ausfd}
\FunctionTok{sudo}\NormalTok{ ./open5gs{-}udmd}
\FunctionTok{sudo}\NormalTok{ ./open5gs{-}pcfd}
\FunctionTok{sudo}\NormalTok{ ./open5gs{-}nrfd}
\FunctionTok{sudo}\NormalTok{ ./open5gs{-}bsfd}
\FunctionTok{sudo}\NormalTok{ ./open5gs{-}hssd}
\end{Highlighting}
\end{Shaded}

Ou tout-en-un :

\begin{Shaded}
\begin{Highlighting}[]
\FunctionTok{sudo}\NormalTok{ ./open5gs{-}daemon}
\end{Highlighting}
\end{Shaded}

\begin{center}\rule{0.5\linewidth}{0.5pt}\end{center}

\subsection{Lancer le gNB srsRAN}\label{lancer-le-gnb-srsran}

\begin{Shaded}
\begin{Highlighting}[]
\ExtensionTok{srsran\_gnb} \AttributeTok{{-}{-}cfg}\NormalTok{ srsran\_project/config/gnb.yaml}
\end{Highlighting}
\end{Shaded}

\begin{center}\rule{0.5\linewidth}{0.5pt}\end{center}

\subsection{(Optionnel) Lancer le srsUE (UE 5G SA
logiciel)}\label{optionnel-lancer-le-srsue-ue-5g-sa-logiciel}

\begin{Shaded}
\begin{Highlighting}[]
\ExtensionTok{srsran\_ue} \AttributeTok{{-}{-}cfg}\NormalTok{ srsran\_project/config/ue.yaml}
\end{Highlighting}
\end{Shaded}

Ou utiliser un vrai smartphone 5G + SDR côté réseau.

\begin{center}\rule{0.5\linewidth}{0.5pt}\end{center}

\subsection{Supervision}\label{supervision}

\begin{itemize}
\tightlist
\item
  Utilise le \textbf{WebUI Open5GS} (si installé) pour la gestion des
  SIM et le monitoring.
\item
  Surveille les logs (attach, registration, PDU session, etc.).
\item
  Une fois le UE attaché, fais un \texttt{ping} (depuis l'UE ou via
  l'interface UPF).
\end{itemize}

\begin{center}\rule{0.5\linewidth}{0.5pt}\end{center}

\subsection{Conseils de dépannage}\label{conseils-de-duxe9pannage}

\begin{itemize}
\tightlist
\item
  \textbf{PLMN, clés et IP doivent correspondre} partout.
\item
  Utilise \texttt{127.0.0.1} ou une IP locale réelle selon ton cas.
\item
  Avec du RF\,: vérifie fréquence/bande/SDR.
\item
  En simulation pure (``null device'') tu peux tout tester sans SDR.
\end{itemize}

\begin{center}\rule{0.5\linewidth}{0.5pt}\end{center}

\subsection{Liens utiles}\label{liens-utiles-2}

\begin{itemize}
\tightlist
\item
  \href{https://docs.srsran.com/projects/5g/en/latest/}{Docs srsRAN
  Project 5G}
\item
  \href{https://open5gs.org/open5gs/docs/}{Docs Open5GS}
\item
  \href{https://github.com/open5gs/open5gs/wiki/WebUI}{WebUI Open5GS}
\item
  \href{https://open5gs.org/open5gs/docs/guide-subscription/}{Guide
  abonnés Open5GS}
\end{itemize}

\end{document}
